\documentclass[10pt,aspectratio=169]{beamer}

% ==========================================================
% Configuración general — basada en el estilo de main.tex
% (tema CambridgeUS, paleta verde, natbib, hyperref)
% ==========================================================
\usepackage[utf8]{inputenc}
\usepackage[T1]{fontenc}
\usepackage[spanish]{babel}
\usepackage{graphicx}
\usepackage{booktabs}
\usepackage{array}
\usepackage{tabularx}
\usepackage{amssymb}
\usepackage{tikz}
\usetikzlibrary{arrows.meta,positioning}
\usepackage{natbib}
\usepackage{hyperref}

\usetheme{CambridgeUS}
\usecolortheme{dolphin}

% Paleta de color verde
\definecolor{myNewColorA}{RGB}{27,94,55}
\definecolor{myNewColorB}{RGB}{46,125,80}
\definecolor{myNewColorC}{RGB}{184,218,191}
\setbeamercolor{structure}{fg=myNewColorA}
\setbeamercolor*{palette primary}{bg=myNewColorC,fg=myNewColorA}
\setbeamercolor*{palette secondary}{bg=myNewColorB, fg=white}
\setbeamercolor*{palette tertiary}{bg=myNewColorA, fg=white}
\setbeamercolor*{titlelike}{fg=myNewColorA}
\setbeamercolor*{title}{bg=myNewColorA, fg=white}
\setbeamercolor*{item}{fg=myNewColorA}
\setbeamercolor*{caption name}{fg=myNewColorA}
\usefonttheme{professionalfonts}

\setbeamertemplate{navigation symbols}{}
\setbeamertemplate{footline}[frame number]
\setbeamercolor{block title}{bg=myNewColorC,fg=black}
\setbeamercolor{block body}{bg=myNewColorC!15,fg=black}
\setbeamerfont{footnote}{size=\tiny}

\setbeamerfont{title}{size=\normalsize}
\setbeamerfont{author}{size=\small}
\setbeamerfont{date}{size=\small}
\setbeamerfont{institute}{size=\small}

\title[Arquitectura híbrida SQL--NoSQL]{Diseño de una arquitectura híbrida SQL--NoSQL para la integración y consulta de datos de incendios forestales en el cantón Loja}
\author{\scshape malan mullo marco vinicio \and pujos culque viviana isabel \and carlos guillermo chuncho m.}
\institute[ESPOCH]{Escuela Superior Politécnica de Chimborazo (ESPOCH) \\ Maestría en Estadística con mención en Ciencia de Datos e Inteligencia Artificial}
\date{2026}

\begin{document}

% ==========================================================
% Diapositiva 1: Carátula
% ==========================================================
\begin{frame}
\titlepage
\end{frame}

% ==========================================================
% Diapositiva 2: Problema de integración de datos
% ==========================================================
\begin{frame}{Problema: integración de datos heterogéneos}

\small
\textbf{Cantón Loja:} analizar incendios forestales requiere relacionar detecciones satelitales, precipitación y ubicación espacial.

\vspace{0.25cm}
\begin{columns}[T,onlytextwidth]

\begin{column}{0.48\textwidth}
\begin{block}{Datos heterogéneos}
\small
\begin{itemize}
\setlength\itemsep{5pt}
\item Fuentes satelitales, climáticas y espaciales.
\item Formatos y niveles de detalle diferentes.
\item Loja: fuerte variación del fuego en el gradiente 1549--2757 m s.n.m.
\end{itemize}
\end{block}
\end{column}

\begin{column}{0.48\textwidth}
\begin{block}{Dificultad de integración}
\small
\begin{itemize}
\setlength\itemsep{5pt}
\item Vincular observaciones por lugar y fecha.
\item Modelos generales explican solo 44.7\,\% de esa variación; por zona, hasta 95\,\%.
\item Obtener consultas consistentes y reproducibles.
\end{itemize}
\end{block}
\end{column}

\end{columns}

\vspace{0.25cm}
\begin{block}{Problema de investigación}
\small
La \textbf{heterogeneidad de las fuentes} dificulta integrar y almacenar los datos con trazabilidad, lo que limita las \textbf{consultas conjuntas y reproducibles} sobre incendios forestales en el cantón Loja.
\end{block}

\vspace{0.1cm}
\tiny
\textbf{Fuente:} Balaguer-Beser et al., \textit{Fire Ecology} 22:46 (2026).

\end{frame}

% ==========================================================
% Diapositiva 3: Pregunta de investigación y objetivo general
% ==========================================================
\begin{frame}{Pregunta de investigación y objetivo general}

\vspace*{\fill}
\begin{block}{Pregunta de investigación}
\normalsize
¿Cómo diseñar una arquitectura híbrida SQL--NoSQL que permita integrar, almacenar y consultar de forma reproducible datos heterogéneos relacionados con incendios forestales en el cantón Loja?
\end{block}

\vspace*{\fill}
\begin{block}{Objetivo general}
\normalsize
Diseñar una solución inicial de arquitectura híbrida que combine una base de datos SQL y una estructura NoSQL para integrar, organizar y consultar datos espaciales, temporales, satelitales y climáticos relacionados con incendios forestales en el cantón Loja.
\end{block}
\vspace*{\fill}

\footnotesize
\textbf{Nota:} se presenta como un diseño inicial de arquitectura de datos, no como un modelo predictivo terminado.

\end{frame}

% ==========================================================
% Diapositiva 4: Inventario de fuentes y tabla de variables
% ==========================================================
\begin{frame}{Inventario de fuentes y variables}

\footnotesize
\textbf{Tabla 1.} Inventario de fuentes de datos --- origen, tipo, forma de captura, formato y almacenamiento previsto.

\scriptsize
\begin{tabularx}{\textwidth}{@{}lXXXX@{}}
\toprule
\textbf{Fuente} & \textbf{Tipo} & \textbf{Captura} & \textbf{Formato} & \textbf{Almacenamiento} \\
\midrule
Malla espacial & Geoespacial estructurado & SIG / Python & GeoJSON, GeoPackage, CSV & PostgreSQL/PostGIS \\
VIIRS & Satelital espacio-temporal & Automática, vía API (Google Earth Engine) & CSV, JSON & MongoDB y PostgreSQL \\
CHIRPS & Climático & Automática, vía API (Google Earth Engine) & CSV, GeoTIFF, Parquet & PostgreSQL y Parquet \\
\bottomrule
\end{tabularx}

\vspace{0.15cm}
\tiny
\textit{Metadatos y riesgos:} se documentan fuente, fecha de captura y resolución; principales riesgos: vacíos temporales, cambios de formato, límites de cuota de la API y desalineación espacial entre fuentes.

\vspace{0.25cm}
\footnotesize
\textbf{Tabla 2.} Diccionario de variables --- nombre, descripción, fuente de origen, tipo de dato y almacenamiento.

\scriptsize
\begin{tabularx}{\textwidth}{@{}lXccX@{}}
\toprule
\textbf{Variable} & \textbf{Descripción} & \textbf{Fuente} & \textbf{Tipo} & \textbf{Almacenamiento} \\
\midrule
\texttt{celda\_id} & Identificador de la celda espacial & Malla & Categórica & PostgreSQL \\
\texttt{latitud/longitud} & Coordenadas geográficas & Malla & Numérica & PostgreSQL/PostGIS \\
\texttt{fecha} & Fecha de observación & VIIRS/CHIRPS & Temporal & PostgreSQL/MongoDB \\
\texttt{frp\_mw} & Potencia radiativa del fuego & VIIRS & Numérica & PostgreSQL/MongoDB \\
\texttt{precipitacion\_*\_mm} & Precipitación acum./media/máx. & CHIRPS & Numérica & PostgreSQL/Parquet \\
\bottomrule
\end{tabularx}

\vspace{0.15cm}
\tiny
\textbf{Unidad de análisis:} celda espacial de 500 m $\times$ 500 m observada durante un mes. \quad
\textbf{Clave lógica:} \texttt{celda\_id + fecha}

\end{frame}

% ==========================================================
% Diapositiva 5: Diseño comparativo SQL--NoSQL
% ==========================================================
\begin{frame}{Diseño comparativo SQL--NoSQL y solución propuesta}

\begin{columns}[T,onlytextwidth]

\begin{column}{0.40\textwidth}
\centering
\begin{tikzpicture}[
  box/.style={draw=myNewColorA, rounded corners, fill=myNewColorC!35,
              align=center, text width=3.2cm, minimum height=0.62cm, font=\tiny, inner sep=1pt},
  arr/.style={-{Stealth[length=1.5mm]}, thick, myNewColorA}
]
\node[box] (a) {Fuentes: CSV/JSON,\\ satelitales y climáticas};
\node[box, below=0.14cm of a] (b) {PostgreSQL/PostGIS\\ + MongoDB};
\node[box, below=0.14cm of b] (c) {Consultas SQL y\\ agregaciones NoSQL};
\node[box, below=0.14cm of c] (d) {Dataset reproducible\\ para análisis};
\draw[arr] (a) -- (b);
\draw[arr] (b) -- (c);
\draw[arr] (c) -- (d);
\end{tikzpicture}

\vspace{0.15cm}
\raggedright
\tiny
\textbf{Entidades principales:}
\begin{itemize}
\setlength\itemsep{0pt}
\item \texttt{dim\_celda}, \texttt{dim\_fecha}
\item \texttt{fact\_incendio}, \texttt{fact\_clima}
\item Colección \texttt{detecciones\_viirs}
\end{itemize}
\end{column}

\begin{column}{0.60\textwidth}
\scriptsize
\textbf{PostgreSQL/PostGIS}
\begin{itemize}
\setlength\itemsep{0pt}
\item Tablas estructuradas
\item Claves primarias y foráneas
\item Relaciones entre celdas, fechas, incendios y clima
\item Consultas \texttt{JOIN}, \texttt{GROUP BY} y agregaciones
\item Almacenamiento de geometrías
\end{itemize}

\vspace{0.2cm}
\textbf{MongoDB}
\begin{itemize}
\setlength\itemsep{0pt}
\item Documentos JSON
\item Atributos flexibles
\item Metadatos e información de calidad
\item Agregaciones documentales
\end{itemize}
\end{column}

\end{columns}

\vspace{0.1cm}
\begin{block}{\scriptsize Pregunta analítica de ejemplo}
\scriptsize ¿Qué celdas y meses presentan detecciones de incendio, mayor FRP y menor precipitación?
\end{block}

\vspace{-0.1cm}
\scriptsize
Solución local, reproducible y que no requiere un servidor dedicado para esta actividad.

\end{frame}

\end{document}
